\documentclass[12pt,a4paper]{article}

\usepackage[T2A]{fontenc}
\usepackage[utf8]{inputenc}
\usepackage[russian]{babel}
\usepackage{amsmath,amssymb,amsthm,mathtools}
\usepackage{array}
\usepackage{booktabs}
\usepackage{algorithm}
\usepackage{algpseudocode}
\usepackage{geometry}
\usepackage{enumitem}

\geometry{left=25mm,right=15mm,top=20mm,bottom=20mm}

\newtheorem{definition}{Определение}
\newtheorem{theorem}{Теорема}
\newtheorem{lemma}{Лемма}
\newtheorem{proposition}{Утверждение}
\newtheorem{example}{Пример}
\newtheorem{corollary}{Заклчение}

\DeclareMathOperator{\sgn}{sgn}
\DeclareMathOperator{\seq}{seq}
\DeclareMathOperator{\ord}{ord}

\title{Билет №19 \\ \small Секвентный анализ. Синтез базисных функций для секвентных преобразований. Базис функций Адамара. Алгоритм быстрого преобразования Адамара-Уолша. (СпБПУ)}
\author{Сериков Владислав Александрович}
\date{5 мая 2026}

\begin{document}

\maketitle
\tableofcontents
\newpage

\section{Идея секвентного анализа}

Секвентный анализ является аналогом гармонического анализа, но вместо синусоидальных базисных функций используются кусочно-постоянные прямоугольные функции, принимающие значения $+1$ и $-1$.

В классическом спектральном анализе сигнал раскладывается по базису комплексных экспонент:
\[
    x(t)=\sum_{k} c_k e^{i2\pi kt/T}.
\]
Число $k$ интерпретируется как частота. В секвентном анализе роль частоты играет \textbf{секвента}.

\begin{definition}[Секвента]
Секвентой функции, заданной на интервале, называется число смен знака функции на этом интервале.
\end{definition}

Например, функция
\[
    (+,+,-,-)
\]
имеет одну смену знака, а функция
\[
    (+,-,+,-)
\]
имеет три смены знака.

Секвентные преобразования особенно удобны для дискретных сигналов, двоичных систем, цифровой обработки изображений, сжатия данных и анализа булевых функций, поскольку базисные функции имеют только два уровня:
\[
    +1,\quad -1.
\]

\section{Дискретное пространство сигналов}

Пусть рассматривается дискретный сигнал длины
\[
    N=2^m,\qquad m\in\mathbb{N}.
\]
Сигнал представляется вектором
\[
    x=
    \begin{pmatrix}
        x_0\\x_1\\ \vdots\\ x_{N-1}
    \end{pmatrix}
    \in \mathbb{R}^N.
\]

Секвентное преобразование ставит в соответствие вектору $x$ набор коэффициентов
\[
    X=
    \begin{pmatrix}
        X_0\\X_1\\ \vdots\\ X_{N-1}
    \end{pmatrix},
\]
где каждый коэффициент показывает степень ``похожести'' сигнала на соответствующую базисную функцию.

Общая форма ортогонального дискретного преобразования:
\[
    X = A x,
\]
где $A$ --- матрица преобразования. Если строки $A$ образуют ортогональный базис, то обратное преобразование имеет вид
\[
    x = A^{-1}X.
\]

Если матрица нормирована так, что
\[
    A A^T = I,
\]
то
\[
    A^{-1}=A^T.
\]

\section{Матрицы Адамара}

\begin{definition}[Матрица Адамара]
Матрицей Адамара порядка $N$ называется квадратная матрица
\[
    H_N\in\{+1,-1\}^{N\times N},
\]
удовлетворяющая условию ортогональности строк:
\[
    H_N H_N^T = N I_N.
\]
\end{definition}

Иначе говоря, любые две различные строки матрицы Адамара ортогональны:
\[
    \sum_{j=0}^{N-1} h_{ij}h_{kj}=0,\qquad i\ne k,
\]
а скалярное произведение строки самой на себя равно $N$:
\[
    \sum_{j=0}^{N-1} h_{ij}^2=N.
\]

\subsection{Матрицы Сильвестра}

Наиболее важный для цифровой обработки случай --- матрицы Адамара порядка $2^m$, строящиеся рекурсивно.

\begin{definition}[Конструкция Сильвестра]
Матрицы Адамара Сильвестра определяются рекуррентно:
\[
    H_1 = [1],
\]
\[
    H_{2N}
    =
    \begin{pmatrix}
        H_N & H_N\\
        H_N & -H_N
    \end{pmatrix}.
\]
Эквивалентно:
\[
    H_{2^m}
    =
    H_2^{\otimes m},
    \qquad
    H_2=
    \begin{pmatrix}
        1 & 1\\
        1 & -1
    \end{pmatrix},
\]
где $\otimes$ обозначает кронекерово произведение.
\end{definition}

Первые матрицы:
\[
H_1=[1],
\]
\[
H_2=
\begin{pmatrix}
1&1\\
1&-1
\end{pmatrix},
\]
\[
H_4=
\begin{pmatrix}
1&1&1&1\\
1&-1&1&-1\\
1&1&-1&-1\\
1&-1&-1&1
\end{pmatrix}.
\]

\begin{theorem}[Ортогональность матриц Сильвестра]
Для любого $m\ge 0$ матрица $H_{2^m}$, построенная по правилу Сильвестра, является матрицей Адамара, то есть
\[
    H_{2^m}H_{2^m}^T = 2^m I_{2^m}.
\]
\end{theorem}

\begin{proof}
Докажем утверждение индукцией по $m$.

\textbf{База индукции.} При $m=0$ имеем
\[
    H_1=[1].
\]
Тогда
\[
    H_1H_1^T=[1]=1\cdot I_1.
\]
Следовательно, утверждение верно.

\textbf{Индукционный переход.} Предположим, что для некоторого $N=2^m$ выполнено
\[
    H_NH_N^T = N I_N.
\]
Рассмотрим матрицу
\[
    H_{2N}=
    \begin{pmatrix}
        H_N&H_N\\
        H_N&-H_N
    \end{pmatrix}.
\]
Тогда
\[
H_{2N}H_{2N}^T
=
\begin{pmatrix}
H_N&H_N\\
H_N&-H_N
\end{pmatrix}
\begin{pmatrix}
H_N^T&H_N^T\\
H_N^T&-H_N^T
\end{pmatrix}.
\]
Выполним блочное умножение:
\[
H_{2N}H_{2N}^T
=
\begin{pmatrix}
H_NH_N^T+H_NH_N^T
&
H_NH_N^T-H_NH_N^T
\\
H_NH_N^T-H_NH_N^T
&
H_NH_N^T+H_NH_N^T
\end{pmatrix}.
\]
То есть
\[
H_{2N}H_{2N}^T
=
\begin{pmatrix}
2H_NH_N^T&0\\
0&2H_NH_N^T
\end{pmatrix}.
\]
По предположению индукции:
\[
    H_NH_N^T = NI_N.
\]
Следовательно,
\[
H_{2N}H_{2N}^T
=
\begin{pmatrix}
2NI_N&0\\
0&2NI_N
\end{pmatrix}
=
2N I_{2N}.
\]
Значит, $H_{2N}$ является матрицей Адамара. Теорема доказана.
\end{proof}

\section{Преобразование Адамара}

\begin{definition}[Ненормированное преобразование Адамара]
Пусть $x\in\mathbb{R}^N$, где $N=2^m$. Ненормированным преобразованием Адамара называется преобразование
\[
    X = H_N x.
\]
\end{definition}

Компоненты вектора $X$:
\[
    X_k = \sum_{n=0}^{N-1} h_{kn}x_n,
    \qquad k=0,1,\ldots,N-1.
\]

\begin{definition}[Нормированное преобразование Адамара]
Нормированное преобразование Адамара задается матрицей
\[
    \widehat{H}_N=\frac{1}{\sqrt{N}}H_N.
\]
Тогда
\[
    X=\widehat{H}_N x.
\]
\end{definition}

\begin{theorem}[Обратимость преобразования Адамара]
Для ненормированного преобразования Адамара
\[
    X=H_Nx
\]
обратное преобразование имеет вид
\[
    x=\frac{1}{N}H_N^T X.
\]
Если $H_N$ симметрична, как матрицы Сильвестра, то
\[
    x=\frac{1}{N}H_NX.
\]
Для нормированной матрицы
\[
    \widehat{H}_N=\frac{1}{\sqrt{N}}H_N
\]
выполнено
\[
    \widehat{H}_N^{-1}=\widehat{H}_N^T.
\]
\end{theorem}

\begin{proof}
Из свойства матрицы Адамара:
\[
    H_NH_N^T=NI_N.
\]
Следовательно,
\[
    H_N^{-1}=\frac{1}{N}H_N^T.
\]
Если
\[
    X=H_Nx,
\]
то, умножая обе части слева на $H_N^{-1}$, получаем
\[
    x=H_N^{-1}X=\frac{1}{N}H_N^TX.
\]
Для нормированной матрицы:
\[
    \widehat{H}_N\widehat{H}_N^T
    =
    \frac{1}{\sqrt{N}}H_N
    \frac{1}{\sqrt{N}}H_N^T
    =
    \frac{1}{N}H_NH_N^T
    =
    I_N.
\]
Следовательно,
\[
    \widehat{H}_N^{-1}=\widehat{H}_N^T.
\]
Теорема доказана.
\end{proof}

\section{Базис функций Адамара--Уолша}

\subsection{Двоичное представление индексов}

Пусть
\[
    N=2^m.
\]
Каждый индекс $n\in\{0,1,\ldots,N-1\}$ можно представить в двоичной форме:
\[
    n = n_{m-1}2^{m-1}+\cdots+n_1 2+n_0,
\]
где
\[
    n_i\in\{0,1\}.
\]
Будем записывать
\[
    n=(n_{m-1},\ldots,n_1,n_0)_2.
\]

Аналогично:
\[
    k=(k_{m-1},\ldots,k_1,k_0)_2.
\]

\subsection{Адамарова, или натуральная, нумерация}

Элементы матрицы Сильвестра можно записать формулой:
\[
    h_{kn}=(-1)^{\langle k,n\rangle},
\]
где
\[
    \langle k,n\rangle
    =
    k_0n_0\oplus k_1n_1\oplus\cdots\oplus k_{m-1}n_{m-1},
\]
а $\oplus$ --- сложение по модулю $2$.

Иначе:
\[
    h_{kn}=(-1)^{\sum_{j=0}^{m-1} k_j n_j}.
\]

Так как важна только четность суммы, запись со сложением по модулю $2$ эквивалентна.

\begin{definition}[Функции Адамара]
Дискретные функции
\[
    \operatorname{had}_k(n)=(-1)^{\langle k,n\rangle},
    \qquad k,n=0,\ldots,N-1,
\]
называются функциями Адамара в натуральном порядке.
\end{definition}

Строки матрицы $H_N$ являются значениями функций Адамара:
\[
    H_N=
    \begin{pmatrix}
        \operatorname{had}_0(0)&\cdots&\operatorname{had}_0(N-1)\\
        \vdots&\ddots&\vdots\\
        \operatorname{had}_{N-1}(0)&\cdots&\operatorname{had}_{N-1}(N-1)
    \end{pmatrix}.
\]

\subsection{Функции Радемахера}

Для синтеза базисных функций удобно использовать функции Радемахера.

\begin{definition}[Дискретные функции Радемахера]
Для $j=0,1,\ldots,m-1$ функция Радемахера определяется как
\[
    r_j(n)=(-1)^{n_j},
\]
где $n_j$ --- $j$-й бит двоичного представления числа $n$.
\end{definition}

Тогда функция Адамара может быть записана как произведение функций Радемахера:
\[
    \operatorname{had}_k(n)
    =
    \prod_{j=0}^{m-1} r_j(n)^{k_j}.
\]

Действительно:
\[
    \prod_{j=0}^{m-1} r_j(n)^{k_j}
    =
    \prod_{j=0}^{m-1} \left((-1)^{n_j}\right)^{k_j}
    =
    (-1)^{\sum_{j=0}^{m-1}k_jn_j}.
\]

\section{Секвентный порядок и функции Уолша}

Матрица Сильвестра задает функции в так называемом \textbf{натуральном}, или \textbf{адамаровом}, порядке. Однако для секвентного анализа удобнее располагать базисные функции по возрастанию числа смен знака.

Такой порядок называется \textbf{секвентным порядком}, а соответствующие функции обычно называют функциями Уолша.

\begin{definition}[Функции Уолша]
Функции Уолша --- это функции Адамара, упорядоченные по возрастанию секвенты, то есть по числу смен знака.
\end{definition}

\subsection{Код Грея и связь порядков}

Переход от секвентного номера $s$ к номеру строки в натуральном адамаровом порядке осуществляется с помощью кода Грея:
\[
    k = g(s)=s\oplus \left\lfloor \frac{s}{2}\right\rfloor.
\]
Здесь $\oplus$ обозначает побитовое исключающее ИЛИ.

Тогда функция Уолша с секвентным номером $s$ определяется как
\[
    \operatorname{wal}_s(n)
    =
    \operatorname{had}_{g(s)}(n).
\]

В двоичной форме код Грея задается так. Если
\[
    s=(s_{m-1},s_{m-2},\ldots,s_0)_2,
\]
то
\[
    g(s)=(g_{m-1},g_{m-2},\ldots,g_0)_2,
\]
где
\[
    g_{m-1}=s_{m-1},
\]
\[
    g_j=s_{j+1}\oplus s_j,\qquad j=0,1,\ldots,m-2.
\]

\subsection{Пример для $N=4$}

Для $N=4$ матрица Адамара в натуральном порядке:
\[
H_4=
\begin{pmatrix}
1&1&1&1\\
1&-1&1&-1\\
1&1&-1&-1\\
1&-1&-1&1
\end{pmatrix}.
\]

Посчитаем секвенты строк:

\[
\begin{array}{c|c|c}
k & \operatorname{had}_k & \text{число смен знака}\\
\hline
0 & (+,+,+,+) & 0\\
1 & (+,-,+,-) & 3\\
2 & (+,+,-,-) & 1\\
3 & (+,-,-,+) & 2
\end{array}
\]

Следовательно, секвентный порядок строк:
\[
    0,\;2,\;3,\;1.
\]

Это соответствует коду Грея:
\[
\begin{array}{c|c|c}
s & g(s)=s\oplus(s\gg 1) & \operatorname{wal}_s\\
\hline
0 & 0 & \operatorname{had}_0\\
1 & 1\oplus 0=1 & \operatorname{had}_1 \quad \text{в некоторых соглашениях}\\
\end{array}
\]

Следует отметить, что существуют разные соглашения о порядке битов и о том, по какому индексу выполняется перестановка. В инженерной литературе часто используется секвентная матрица Уолша:
\[
W_4=
\begin{pmatrix}
1&1&1&1\\
1&1&-1&-1\\
1&-1&-1&1\\
1&-1&1&-1
\end{pmatrix}.
\]
Ее строки имеют соответственно $0,1,2,3$ смены знака.

\section{Синтез базисных функций для секвентных преобразований}

\subsection{Задача синтеза}

Синтез базисных функций означает построение набора функций
\[
    \varphi_0,\varphi_1,\ldots,\varphi_{N-1},
\]
которые удовлетворяют следующим условиям:

\begin{enumerate}
    \item принимают значения $+1$ и $-1$;
    \item попарно ортогональны:
    \[
        \sum_{n=0}^{N-1}\varphi_i(n)\varphi_j(n)=0,\qquad i\ne j;
    \]
    \item имеют одинаковую энергию:
    \[
        \sum_{n=0}^{N-1}\varphi_i^2(n)=N;
    \]
    \item упорядочены по возрастанию секвенты.
\end{enumerate}

\subsection{Метод 1: рекурсивный синтез через матрицы Сильвестра}

Алгоритм:

\begin{enumerate}
    \item Начать с
    \[
        H_1=[1].
    \]
    \item На каждом шаге строить
    \[
        H_{2N}=
        \begin{pmatrix}
            H_N&H_N\\
            H_N&-H_N
        \end{pmatrix}.
    \]
    \item Полученные строки являются функциями Адамара.
    \item Для получения секвентного порядка выполнить перестановку строк по числу смен знака.
\end{enumerate}

\subsection{Метод 2: синтез через функции Радемахера}

Пусть $N=2^m$. Для каждого $j=0,\ldots,m-1$ строится функция
\[
    r_j(n)=(-1)^{n_j}.
\]
Далее для каждого номера $k$ строится функция
\[
    \operatorname{had}_k(n)=\prod_{j=0}^{m-1}r_j(n)^{k_j}.
\]

Если нужен секвентный номер $s$, то сначала находится соответствующий номер $k$ через код Грея или обратный код Грея, после чего строится
\[
    \operatorname{wal}_s(n)=\operatorname{had}_k(n).
\]

\subsection{Пример синтеза для $N=8$}

Пусть $N=8$, тогда $m=3$, а $n$ принимает значения $0,\ldots,7$.

Двоичные представления:
\[
\begin{array}{c|c}
n & (n_2n_1n_0)_2\\
\hline
0&000\\
1&001\\
2&010\\
3&011\\
4&100\\
5&101\\
6&110\\
7&111
\end{array}
\]

Функции Радемахера:
\[
r_0(n)=(-1)^{n_0}=(+,-,+,-,+,-,+,-),
\]
\[
r_1(n)=(-1)^{n_1}=(+,+,-,-,+,+,-,-),
\]
\[
r_2(n)=(-1)^{n_2}=(+,+,+,+,-,-,-,-).
\]

Например, для $k=5=(101)_2$:
\[
    \operatorname{had}_5(n)
    =
    r_2(n)r_0(n).
\]
Следовательно,
\[
\operatorname{had}_5(n)
=
(+,-,+,-,-,+,-,+).
\]

\section{Ортогональность базиса Адамара--Уолша}

\begin{theorem}[Ортогональность функций Адамара]
Пусть
\[
    \operatorname{had}_k(n)=(-1)^{\langle k,n\rangle},
    \qquad k,n\in\{0,\ldots,2^m-1\}.
\]
Тогда
\[
    \sum_{n=0}^{N-1}\operatorname{had}_k(n)\operatorname{had}_l(n)
    =
    \begin{cases}
        N, & k=l,\\
        0, & k\ne l.
    \end{cases}
\]
\end{theorem}

\begin{proof}
Рассмотрим сумму
\[
S=\sum_{n=0}^{N-1}\operatorname{had}_k(n)\operatorname{had}_l(n).
\]
По определению:
\[
\operatorname{had}_k(n)\operatorname{had}_l(n)
=
(-1)^{\langle k,n\rangle}
(-1)^{\langle l,n\rangle}
=
(-1)^{\langle k,n\rangle+\langle l,n\rangle}.
\]
Так как показатели рассматриваются по модулю $2$,
\[
    \langle k,n\rangle+\langle l,n\rangle
    =
    \langle k\oplus l,n\rangle.
\]
Следовательно,
\[
    S=
    \sum_{n=0}^{N-1}(-1)^{\langle k\oplus l,n\rangle}.
\]

Если $k=l$, то $k\oplus l=0$, поэтому
\[
    \langle k\oplus l,n\rangle=0
\]
для всех $n$, и
\[
    S=\sum_{n=0}^{N-1}1=N.
\]

Пусть теперь $k\ne l$. Тогда
\[
    a=k\oplus l\ne 0.
\]
Значит, существует хотя бы один бит $a_j=1$. Разобьем все векторы $n$ на пары, отличающиеся только в $j$-м бите. Если $n'$ получается из $n$ изменением $j$-го бита, то
\[
    \langle a,n'\rangle
    =
    \langle a,n\rangle\oplus 1.
\]
Следовательно,
\[
    (-1)^{\langle a,n'\rangle}
    =
    -(-1)^{\langle a,n\rangle}.
\]
Сумма по каждой такой паре равна нулю. Все $2^m$ векторов разбиваются на такие пары, значит
\[
    S=0.
\]
Теорема доказана.
\end{proof}

\begin{corollary}
Система нормированных функций
\[
    \psi_k(n)=\frac{1}{\sqrt{N}}\operatorname{had}_k(n)
\]
является ортонормированным базисом в $\mathbb{R}^N$.
\end{corollary}

\begin{proof}
Из предыдущей теоремы следует ортонормированность:
\[
    \sum_{n=0}^{N-1}\psi_k(n)\psi_l(n)
    =
    \frac{1}{N}
    \sum_{n=0}^{N-1}\operatorname{had}_k(n)\operatorname{had}_l(n)
    =
    \delta_{kl}.
\]
Так как имеется ровно $N$ попарно ортонормированных векторов в $N$-мерном пространстве, они образуют базис.
\end{proof}

\section{Преобразование Уолша--Адамара}

\begin{definition}[Преобразование Уолша--Адамара]
Пусть $W_N$ --- матрица, строки которой являются функциями Уолша. Тогда преобразованием Уолша--Адамара называется
\[
    X = W_N x.
\]
\end{definition}

Если используется нормированная матрица
\[
    \widehat{W}_N=\frac{1}{\sqrt{N}}W_N,
\]
то
\[
    X=\widehat{W}_N x,
\qquad
    x=\widehat{W}_N^T X.
\]

В ненормированном случае:
\[
    x=\frac{1}{N}W_N^T X.
\]

Если $W_N$ получается из $H_N$ перестановкой строк, то
\[
    W_N = P H_N,
\]
где $P$ --- перестановочная матрица. Поскольку
\[
    P^TP=I,
\]
ортогональность сохраняется:
\[
    W_NW_N^T
    =
    P H_N H_N^T P^T
    =
    P(NI)P^T
    =
    N I.
\]

\section{Быстрое преобразование Адамара--Уолша}

\subsection{Необходимость быстрого алгоритма}

Прямое вычисление
\[
    X=H_Nx
\]
требует порядка
\[
    N^2
\]
операций сложения и вычитания.

Для больших $N$ это неэффективно. Благодаря рекурсивной структуре матриц Адамара можно вычислять преобразование за
\[
    O(N\log_2 N)
\]
операций.

\subsection{Базовая идея}

Пусть
\[
    x=
    \begin{pmatrix}
        x^{(0)}\\
        x^{(1)}
    \end{pmatrix},
\]
где $x^{(0)}$ и $x^{(1)}$ --- первая и вторая половины вектора длины $N/2$.

Тогда
\[
H_N=
\begin{pmatrix}
H_{N/2}&H_{N/2}\\
H_{N/2}&-H_{N/2}
\end{pmatrix}.
\]

Следовательно,
\[
H_Nx
=
\begin{pmatrix}
H_{N/2}&H_{N/2}\\
H_{N/2}&-H_{N/2}
\end{pmatrix}
\begin{pmatrix}
x^{(0)}\\
x^{(1)}
\end{pmatrix}
=
\begin{pmatrix}
H_{N/2}(x^{(0)}+x^{(1)})\\
H_{N/2}(x^{(0)}-x^{(1)})
\end{pmatrix}.
\]

Таким образом, задача размера $N$ сводится к двум задачам размера $N/2$:
\[
    H_Nx
    =
    \begin{pmatrix}
        H_{N/2}(x^{(0)}+x^{(1)})\\
        H_{N/2}(x^{(0)}-x^{(1)})
    \end{pmatrix}.
\]

\subsection{Операция ``бабочка''}

Основная операция быстрого преобразования:
\[
    a'=a+b,
\]
\[
    b'=a-b.
\]

Графически:
\[
\begin{array}{ccc}
a & \longrightarrow & a+b\\
b & \longrightarrow & a-b
\end{array}
\]

Эта операция называется \textbf{бабочкой}.

\section{Алгоритм быстрого преобразования Адамара--Уолша}

\subsection{Итеративный алгоритм для натурального порядка}

Пусть задан вектор
\[
    x=(x_0,x_1,\ldots,x_{N-1}),
\qquad N=2^m.
\]
Требуется вычислить
\[
    X=H_Nx.
\]

\begin{algorithm}[H]
\caption{Быстрое преобразование Адамара--Уолша}
\begin{algorithmic}[1]
\Require Вектор $x$ длины $N=2^m$
\Ensure Вектор $X=H_Nx$
\State $X \gets x$
\State $h \gets 1$
\While{$h<N$}
    \For{$i=0$ to $N-1$ step $2h$}
        \For{$j=i$ to $i+h-1$}
            \State $a \gets X_j$
            \State $b \gets X_{j+h}$
            \State $X_j \gets a+b$
            \State $X_{j+h} \gets a-b$
        \EndFor
    \EndFor
    \State $h \gets 2h$
\EndWhile
\State \Return $X$
\end{algorithmic}
\end{algorithm}

После выполнения алгоритма получается ненормированное преобразование. Для нормированного преобразования необходимо дополнительно выполнить:
\[
    X\gets \frac{1}{\sqrt{N}}X.
\]

\subsection{Пояснение шагов алгоритма}

На первом шаге объединяются соседние элементы:
\[
(x_0,x_1),\quad (x_2,x_3),\quad \ldots
\]
и каждая пара заменяется на
\[
(x_{2i}+x_{2i+1},\; x_{2i}-x_{2i+1}).
\]

На втором шаге объединяются блоки длины $4$:
\[
(x_0,x_1,x_2,x_3),
\]
и выполняются операции:
\[
(x_0,x_2),\quad (x_1,x_3).
\]

На третьем шаге объединяются блоки длины $8$ и так далее.

\subsection{Минимальный пример для $N=4$}

Пусть
\[
    x=(1,2,3,4)^T.
\]

Матрица:
\[
H_4=
\begin{pmatrix}
1&1&1&1\\
1&-1&1&-1\\
1&1&-1&-1\\
1&-1&-1&1
\end{pmatrix}.
\]

Прямое вычисление:
\[
X=H_4x=
\begin{pmatrix}
1+2+3+4\\
1-2+3-4\\
1+2-3-4\\
1-2-3+4
\end{pmatrix}
=
\begin{pmatrix}
10\\
-2\\
-4\\
0
\end{pmatrix}.
\]

Теперь выполним быстрый алгоритм.

Исходный вектор:
\[
    (1,2,3,4).
\]

\textbf{Шаг 1.} Размер бабочки $h=1$.

Пары:
\[
(1,2)\mapsto (1+2,1-2)=(3,-1),
\]
\[
(3,4)\mapsto (3+4,3-4)=(7,-1).
\]
Получаем:
\[
    (3,-1,7,-1).
\]

\textbf{Шаг 2.} Размер бабочки $h=2$.

Пары внутри блока длины $4$:
\[
(3,7)\mapsto(3+7,3-7)=(10,-4),
\]
\[
(-1,-1)\mapsto(-1+(-1),-1-(-1))=(-2,0).
\]

С учетом позиций:
\[
    (10,-2,-4,0).
\]

Следовательно,
\[
    X=(10,-2,-4,0)^T,
\]
что совпадает с прямым умножением на $H_4$.

\subsection{Пример для $N=8$}

Пусть
\[
    x=(1,1,1,1,0,0,0,0).
\]

\textbf{Шаг 1:} $h=1$.

\[
(1,1)\mapsto(2,0),
\quad
(1,1)\mapsto(2,0),
\quad
(0,0)\mapsto(0,0),
\quad
(0,0)\mapsto(0,0).
\]
Получаем:
\[
    (2,0,2,0,0,0,0,0).
\]

\textbf{Шаг 2:} $h=2$.

\[
(2,2)\mapsto(4,0),
\quad
(0,0)\mapsto(0,0),
\]
\[
(0,0)\mapsto(0,0),
\quad
(0,0)\mapsto(0,0).
\]
Получаем:
\[
    (4,0,0,0,0,0,0,0).
\]

\textbf{Шаг 3:} $h=4$.

\[
(4,0)\mapsto(4,4),
\]
\[
(0,0)\mapsto(0,0),
\]
\[
(0,0)\mapsto(0,0),
\]
\[
(0,0)\mapsto(0,0).
\]
Получаем:
\[
    X=(4,0,0,0,4,0,0,0).
\]

\section{Доказательство корректности быстрого преобразования}

\begin{theorem}[Корректность быстрого преобразования Адамара]
Итеративный алгоритм быстрого преобразования Адамара вычисляет
\[
    X=H_Nx
\]
для любого вектора $x\in\mathbb{R}^N$, где $N=2^m$.
\end{theorem}

\begin{proof}
Докажем утверждение по индукции по $m$, где $N=2^m$.

\textbf{База.} При $N=1$ матрица
\[
    H_1=[1],
\]
и преобразование тождественно:
\[
    X=x.
\]
Алгоритм не выполняет ни одной операции и возвращает $x$, то есть результат верен.

\textbf{Индукционный переход.} Предположим, что алгоритм корректно вычисляет преобразование для длины $N/2$.

Пусть
\[
    x=
    \begin{pmatrix}
        x^{(0)}\\
        x^{(1)}
    \end{pmatrix},
\]
где $x^{(0)},x^{(1)}\in\mathbb{R}^{N/2}$.

По определению матрицы Сильвестра:
\[
    H_N=
    \begin{pmatrix}
        H_{N/2}&H_{N/2}\\
        H_{N/2}&-H_{N/2}
    \end{pmatrix}.
\]
Тогда
\[
H_Nx
=
\begin{pmatrix}
H_{N/2}(x^{(0)}+x^{(1)})\\
H_{N/2}(x^{(0)}-x^{(1)})
\end{pmatrix}.
\]

Алгоритм на соответствующем уровне выполняет операции
\[
    x^{(0)}+x^{(1)}
\]
и
\[
    x^{(0)}-x^{(1)}
\]
покомпонентно с помощью бабочек. Затем к каждой половине применяется тот же алгоритм для размера $N/2$.

По предположению индукции, каждая из этих двух половин преобразуется корректно:
\[
    x^{(0)}+x^{(1)}
    \mapsto
    H_{N/2}(x^{(0)}+x^{(1)}),
\]
\[
    x^{(0)}-x^{(1)}
    \mapsto
    H_{N/2}(x^{(0)}-x^{(1)}).
\]
Следовательно, весь результат равен
\[
    H_Nx.
\]
Теорема доказана.
\end{proof}

\section{Оценка вычислительной сложности}

На каждом уровне алгоритма выполняется $N/2$ бабочек. Каждая бабочка требует:
\[
    1 \text{ сложение} + 1 \text{ вычитание}.
\]
Всего уровней:
\[
    \log_2 N.
\]
Поэтому число бабочек:
\[
    \frac{N}{2}\log_2 N.
\]
Число арифметических операций:
\[
    N\log_2 N.
\]

Таким образом:
\[
    T(N)=O(N\log N).
\]

Для сравнения, прямое умножение матрицы $H_N$ на вектор требует
\[
    O(N^2)
\]
операций.

\section{Секвентно упорядоченное быстрое преобразование}

Быстрое преобразование в стандартной форме обычно выдает коэффициенты в адамаровом, или натуральном, порядке. Если требуется секвентный порядок, применяется перестановка индексов.

Пусть
\[
    X^{(H)}=H_Nx
\]
--- коэффициенты в натуральном порядке.

Тогда коэффициенты в секвентном порядке:
\[
    X^{(W)}_s = X^{(H)}_{p(s)},
\]
где $p(s)$ --- перестановка, задающая соответствие между секвентным и адамаровым порядком.

Один из распространенных способов:
\[
    p(s)=s\oplus(s\gg 1),
\]
где $s\gg 1$ --- сдвиг двоичной записи $s$ вправо на один бит.

При необходимости также может использоваться битовая реверсия, поскольку конкретная форма перестановки зависит от соглашения о порядке битов и от того, используется ли матрица Пэли, Сильвестра или другая форма записи функций Уолша.

\section{Свойства преобразования Адамара--Уолша}

\subsection{Линейность}

Для любых сигналов $x,y\in\mathbb{R}^N$ и чисел $\alpha,\beta\in\mathbb{R}$:
\[
    H_N(\alpha x+\beta y)
    =
    \alpha H_Nx+\beta H_Ny.
\]

\subsection{Сохранение энергии}

Для нормированного преобразования
\[
    X=\widehat{H}_Nx
\]
выполнено:
\[
    \|X\|_2^2=\|x\|_2^2.
\]

\begin{proof}
Так как
\[
    \widehat{H}_N^T\widehat{H}_N=I,
\]
то
\[
    \|X\|_2^2
    =
    X^TX
    =
    (\widehat{H}_Nx)^T(\widehat{H}_Nx)
    =
    x^T\widehat{H}_N^T\widehat{H}_Nx
    =
    x^Tx
    =
    \|x\|_2^2.
\]
\end{proof}

\subsection{Самообратимость}

Для матриц Сильвестра:
\[
    H_N^T=H_N.
\]
Поэтому
\[
    H_N^{-1}=\frac{1}{N}H_N.
\]
Следовательно, ненормированное преобразование является самообратимым с точностью до множителя $N$:
\[
    H_N(H_Nx)=N x.
\]

Для нормированного преобразования:
\[
    \widehat{H}_N^2=I.
\]

\section{Интерпретация коэффициентов}

Коэффициент преобразования
\[
    X_k=\sum_{n=0}^{N-1}x_n\varphi_k(n)
\]
есть скалярное произведение сигнала $x$ с $k$-й базисной функцией:
\[
    X_k=\langle x,\varphi_k\rangle.
\]

Если $X_k$ велик по модулю, значит сигнал содержит заметную компоненту, похожую на соответствующую функцию Уолша.

Низкие секвенты соответствуют медленно меняющимся прямоугольным функциям. Высокие секвенты соответствуют быстро меняющимся функциям.

\section{Сравнение преобразования Фурье и Уолша--Адамара}

\[
\begin{array}{c|c|c}
\text{Свойство} & \text{Фурье} & \text{Уолш--Адамар}\\
\hline
\text{Базис} & \sin,\cos,e^{i\omega t} & \pm 1\text{-функции}\\
\text{Параметр порядка} & \text{частота} & \text{секвента}\\
\text{Тип функций} & \text{гармонические} & \text{прямоугольные}\\
\text{Операции} & \text{сложения и умножения} & \text{сложения и вычитания}\\
\text{Быстрый алгоритм} & \text{FFT} & \text{FWHT}\\
\text{Сложность} & O(N\log N) & O(N\log N)
\end{array}
\]

\section{Краткая схема применения}

Для анализа сигнала $x$:

\begin{enumerate}
    \item выбрать длину $N=2^m$;
    \item построить или неявно задать базис Адамара--Уолша;
    \item вычислить коэффициенты:
    \[
        X=H_Nx
    \]
    или
    \[
        X=W_Nx;
    \]
    \item при необходимости нормировать:
    \[
        X\gets \frac{1}{\sqrt{N}}X;
    \]
    \item анализировать распределение энергии по секвентам:
    \[
        |X_k|^2.
    \]
\end{enumerate}

Для восстановления сигнала:
\[
    x=\frac{1}{N}H_NX
\]
в ненормированном случае или
\[
    x=\widehat{H}_NX
\]
в нормированном случае.

\section{Основные формулы}

\[
    H_1=[1],
\]
\[
    H_{2N}=
    \begin{pmatrix}
        H_N&H_N\\
        H_N&-H_N
    \end{pmatrix}.
\]

\[
    H_NH_N^T=NI_N.
\]

\[
    X=H_Nx.
\]

\[
    x=\frac{1}{N}H_N^TX.
\]

\[
    \widehat{H}_N=\frac{1}{\sqrt{N}}H_N.
\]

\[
    \operatorname{had}_k(n)=(-1)^{\langle k,n\rangle}.
\]

\[
    \langle k,n\rangle
    =
    \bigoplus_{j=0}^{m-1}k_jn_j.
\]

\[
    \operatorname{had}_k(n)
    =
    \prod_{j=0}^{m-1}r_j(n)^{k_j}.
\]

\[
    r_j(n)=(-1)^{n_j}.
\]

\[
    a'=a+b,\qquad b'=a-b.
\]

\[
    T(N)=O(N\log_2N).
\]

\section{Вывод}

Секвентный анализ основан на разложении сигналов по ортогональному базису прямоугольных функций, принимающих значения $+1$ и $-1$. Основным базисом является базис функций Адамара--Уолша. Матрицы Адамара порядка $2^m$ удобно строятся рекурсивно по конструкции Сильвестра. Их строки образуют ортогональный базис в $\mathbb{R}^{2^m}$.

Функции Адамара в натуральном порядке можно синтезировать через двоичные индексы или через произведения функций Радемахера. Функции Уолша получаются из функций Адамара перестановкой строк в порядке возрастания секвенты.

Быстрое преобразование Адамара--Уолша использует рекурсивную структуру матрицы Сильвестра и операцию бабочки
\[
    (a,b)\mapsto(a+b,a-b).
\]
Это позволяет уменьшить сложность вычислений с
\[
    O(N^2)
\]
до
\[
    O(N\log N).
\]

\end{document}
